
\begin{register}{H}{dcsr}{0x7B0}
\label{regdesc:DCSR}
\regfield{xdebugver}{4}{28}{{4}}%
\regfield{reserved}{12}{16}{{0}}%
\regfield{ebreakm}{1}{15}{{0}}%
\regfield{reserved}{2}{13}{{0}}%
\regfield{ebreaku}{1}{12}{{0}}%
\regfield{stepie}{1}{11}{{0}}%
\regfield{stopcount}{1}{10}{{0}}%
\regfield{stoptime}{1}{9}{{0}}%
\regfield{cause}{3}{6}{{0}}%
\regfield{reserved}{1}{5}{{0}}%
\regfield{mprven}{1}{4}{{0}}%
\regfield{reserved}{1}{3}{{0}}%
\regfield{step}{1}{2}{{0}}%
\regfield{prv}{2}{0}{{0}}%
\reglabel{Reset}\regnewline%
\begin{regdesc}
\begin{reglist}
\label{fielddesc:DCSRXEDBGVER}\item [xdebugver] 表示调试版本。 \\
4：存在外部调试支持\\
(RO)
\label{fielddesc:DCSREBREAKM}\item [ebreakm] 当为 1 时，机器模式下的 ebreak 指令进入调试模式。 \\
(R/W)
\label{fielddesc:DCSREBREAKU}\item [ebreaku] 当为 1 时，用户/程序模式下的 ebreak 指令进入调试模式。 \\
(R/W)
\label{fielddesc:DCSRSTEPIE}\item [stepie] 配置调试模式下的中断控制。\\
0：单步执行时禁用中断 \\
1：单步执行时启用中断 \\
调试器在 hart 运行时不得更改此位的值。\\
(R/W)
\label{fielddesc:DCSRSTOPCOUNT}\item [stopcount] 配置调试模式下的 mcycle 计数器。\\
0：累加计数器\\
1：在调试模式或 ebreak 指令执行期间（并导致进入调试模式），不累加任何计数器，包括 cycle 和 instret CSR\\
(R/W)
\label{fielddesc:DCSRSTOPTIME}\item [stoptime] 配置调试模式下的计时器。\\
0：累加计时器\\
1：在调试模式下不累加任何 hart 本地计时器 \\
(R/W)
\label{fielddesc:DCSRCAUSE}\item [cause] 配置进入调试模式的原因。当单个周期中有多个原因导致进入调试模式，则写入优先级最高的原因。\\
1：执行了 ebreak 指令（优先级 3）\\
2：触发模块导致暂停（优先级 4，最高）\\
3：haltreq 被设置（优先级 1）\\
4：由于 step 被设置，CPU 核心单步执行（优先级 0，最低）\\
5：由于 resethaltreq，hart 在复位后暂停（优先级 2）\\
其他值：保留，供未来使用\\
(RO)
\label{fielddesc:DCSRMPRVEN}\item [mprven] 配置在调试模式下是否使能 mstatus CSR 中的 MPRV 功能。\\
0：禁用 \\
1：使能\\
(R/W)
\item [见下页……]
\end{reglist}\end{regdesc}
\end{register}
\addtocounter{Regfloat}{-1}
\begin{register}{H}{dcsr}{0x7B0}
\begin{regdesc}\begin{reglist}
\item [接上页……]
\label{fielddesc:DCSRSTEP}\item [step] 当被置位且不处于调试模式时，核心将仅执行单个指令，然后进入调试模式。 \\
如果指令由于异常而未能完成，则核心将在执行异常处理程序之前立即进入调试模式，并置位相应的异常寄存器。\\
设置该位不会屏蔽中断，这与 RISC-V 外部调试支持规范 v0.13 中的规定不同。\\
(R/W)
\label{fielddesc:DCSRPRV}\item [prv] 保存核心进入调试模式时的特权级别。退出调试模式时，调试器可以更改此值以改变核心的特权级别。仅支持 \textbf{0x3}（机器模式）和 \textbf{0x0}（用户模式）。(R/W)
\end{reglist}
\end{regdesc}
\end{register}

\begin{register}{H}{dpc}{0x7B1}
\label{regdesc:DPC}
\regfield{dpc}{32}{0}{{0}}%
\reglabel{Reset}\regnewline%
\begin{regdesc}
\begin{reglist}
\label{fielddesc:DPC}\item [dpc] 进入调试模式后，dpc 将写入下一条指令的虚拟地址。恢复执行时，CPU 核心的 PC 将更新为 dpc 保存的虚拟地址。调试器可以写入 dpc 来配置 CPU 恢复执行的位置。(R/W)
\end{reglist}
\end{regdesc}
\end{register}

\begin{longtable}{| l| p{12cm}|}
\caption{进入调试模式时 DPC 中的虚拟地址} \label{tab:riscvcpu-dbg-cause}
\\ \hline
\rowcolor{lightgray}
\textbf{原因} & \textbf{DPC 中的虚拟地址}   \\ \hline
ebreak \label{ebreak} &  ebreak指令的地址  \\ \hline
单步执行 \label{sstep} & 如果没有调试进行，将要执行的下一条指令的地址，即对于不改变程序流的 32 位指令为 pc+4，跳转/分支为目标 PC 等。  \\ \hline
触发模块 \label{trigger} &  如果 \hyperref[fielddesc:MCONTROLTIMING]{timing} 为 0，则为导致触发的指令的地址。如果 \hyperref[fielddesc:MCONTROLTIMING]{timing} 为 1，则为进入调试模式时要执行的下一条指令的地址。\\ \hline
暂停请求\label{halt} & 进入调试模式时要执行的下一条指令的地址。   \\ \hline
\end{longtable}

\begin{register}{H}{dscratch0}{0x7B2}
\label{regdesc:DSCRATCH0}
\regfield{dscratch0}{32}{0}{{0}}%
\reglabel{Reset}\regnewline%
\begin{regdesc}
\begin{reglist}
\label{fielddesc:DSCRATCH0}\item [dscratch0] 由调试模块内部使用。(R/W)
\end{reglist}
\end{regdesc}
\end{register}

\begin{register}{H}{dscratch1}{0x7B3}
\label{regdesc:DSCRATCH1}
\regfield{dscratch1}{32}{0}{{0}}%
\reglabel{Reset}\regnewline%
\begin{regdesc}
\begin{reglist}
\label{fielddesc:DSCRATCH1}\item [dscratch1] 由调试模块内部使用。 (R/W)
\end{reglist}
\end{regdesc}
\end{register}